\chapter{Implementação}
\label{cap:implementacao}

Este capítulo descreve a implementação funcional desenvolvida para materializar a solução proposta. Enquanto o capítulo anterior apresentou a solução em nível conceitual, este capítulo explica como ela foi construída tecnicamente, abordando sua arquitetura, fluxo de simplificação, representação dos perfis de paciente, construção dos prompts, integração com modelos de linguagem, guardrails, recursos auxiliares, indicador de legibilidade e interface de usuário.

A aplicação desenvolvida recebeu o nome de SimplyHealth. Seu objetivo é apoiar a simplificação e adaptação comunicacional de textos técnicos de saúde, considerando diferentes perfis de pacientes e preservando o conteúdo clínico essencial do material original. A solução não atua como sistema clínico autônomo, não gera diagnósticos, não prescreve tratamentos e não valida clinicamente o conteúdo recebido. Sua função é transformar linguisticamente textos previamente definidos ou validados por profissionais de saúde, tornando-os mais compreensíveis para usuários finais.

A implementação foi organizada em dois módulos principais: um backend, desenvolvido em TypeScript com NestJS, e um frontend, desenvolvido com React 18 e Vite. Esses módulos se comunicam por meio de requisições HTTP em formato JSON, seguindo uma arquitetura baseada em API REST.

O backend concentra a lógica principal da aplicação. Ele recebe textos ou arquivos enviados pelo usuário, valida os dados de entrada, identifica o perfil comunicacional selecionado, extrai conteúdo textual quando necessário, constrói os prompts, aciona o provedor de LLM, aplica mecanismos de guardrail e calcula o indicador de legibilidade. O frontend, por sua vez, oferece a interface de interação, permitindo inserir o conteúdo, selecionar perfis, configurar opções de processamento e visualizar os resultados.

A organização interna do backend segue uma divisão em camadas, composta principalmente por controllers, use cases e services. Os controllers recebem as requisições HTTP e realizam validações iniciais. Os use cases coordenam o fluxo da aplicação, articulando extração de texto, construção do prompt, chamada ao LLM, aplicação dos guardrails e montagem da resposta. Os services encapsulam funcionalidades específicas, como comunicação com provedores de LLM, cálculo de legibilidade, processamento de arquivos, avaliação de fidelidade entre o texto original e o texto gerado, além da geração de elementos complementares, como glossário, perguntas frequentes e respostas do chat contextual.

Essa separação evita que a lógica da solução fique concentrada diretamente nos endpoints da API, distribuindo as responsabilidades entre camadas específicas. Com isso, o backend mantém a orquestração do fluxo separada de funcionalidades como extração de texto, chamada ao LLM, cálculo de legibilidade e aplicação de guardrails.

De forma geral, o frontend atua como camada de interação com o usuário, enquanto o backend concentra a orquestração do processamento. A simplificação textual envolve etapas de validação da entrada, extração de conteúdo, construção do prompt, chamada ao LLM, aplicação de guardrails, cálculo de legibilidade e retorno dos resultados à interface. Esse fluxo foi apresentado conceitualmente na Seção~\ref{sec:fluxo-simplificacao-llms}.

As próximas seções detalham os principais elementos da implementação. Inicialmente, são apresentadas as tecnologias utilizadas no backend, no frontend e na integração com os modelos de linguagem. Em seguida, descreve-se a representação dos perfis simulados de pacientes e a forma como esses perfis orientam a adaptação comunicacional. Posteriormente, são discutidas a construção dos prompts, a integração com LLMs, o cálculo do indicador de legibilidade e, por fim, a interface da aplicação desenvolvida.

\section{Tecnologias utilizadas}
\label{sec:tecnologias-utilizadas}

A implementação utilizou tecnologias voltadas ao desenvolvimento web moderno, com separação entre frontend e backend. No backend, foram utilizados TypeScript e NestJS. O TypeScript contribui para a organização do código por meio de tipagem estática, enquanto o NestJS oferece uma estrutura modular baseada em controllers, providers, injeção de dependência e services. Além dessas características técnicas, a escolha também foi influenciada pela experiência prévia do autor com essa stack, o que reduziu o risco de implementação e permitiu concentrar o esforço do trabalho na lógica de simplificação, nos guardrails e na integração com LLMs.

O backend também utiliza recursos para recebimento de arquivos, especialmente no fluxo de upload de PDFs e arquivos TXT. Para a comunicação com modelos de linguagem, a aplicação foi organizada a partir de uma abstração de provedores de LLM. O provedor principal utilizado é o Gemini, da Google. A implementação também prevê suporte ao Claude, da Anthropic, e a um provedor mockado, usado em desenvolvimento e testes para simular respostas sem depender de chamadas externas. Todos os testes e a avaliação descritos neste trabalho foram realizados exclusivamente com o Gemini. 

A escolha do Gemini foi motivada pela facilidade de integração via API, pela disponibilidade de modelos adequados às tarefas de geração e avaliação textual e pelo custo compatível com o escopo de uma implementação acadêmica. Como o objetivo deste trabalho não é treinar ou comparar modelos de linguagem, mas investigar seu uso em um fluxo de simplificação com guardrails, a adoção de um provedor externo permitiu concentrar a implementação na aplicação proposta.

As chamadas à API do Gemini foram configuradas com valores de temperatura distintos conforme o papel. Para a simplificação textual, foi configurada temperatura 0,3. A escolha por um valor reduzido, mas não mínimo, busca equilibrar consistência e fluência: temperaturas muito próximas de zero tendem a produzir textos repetitivos e pouco naturais, enquanto valores altos aumentam a variabilidade de forma indesejada para um contexto clínico. Para os guardrails, foi adotada temperatura 0,1, próxima do determinístico, adequada a tarefas de classificação e avaliação estruturada em que respostas divergentes entre execuções seriam indesejáveis. Os demais parâmetros de geração não foram configurados explicitamente, mantendo os valores padrão da API.

A Tabela~\ref{tab:modelos-llm} apresenta os modelos Gemini utilizados em cada papel da aplicação durante a avaliação.

\begin{table}[htbp]
  \centering
  \caption{Modelos de linguagem utilizados por papel na avaliação da SimplyHealth.}
  \label{tab:modelos-llm}
  \small
  \renewcommand{\arraystretch}{1.15}
  \begin{tabular}{ll}
    \hline
    \textbf{Papel} & \textbf{Modelo} \\
    \hline
    Simplificação textual & gemini-2.5-flash-lite \\
    Transcrição de voz (STT) & gemini-2.5-flash-lite \\
    Síntese de fala (TTS) & gemini-2.5-flash-preview-tts \\
    Guardrail de entrada do chat & gemini-2.5-flash-lite \\
    Guardrail de fidelidade (LLM avaliador) & gemini-2.5-flash \\
    \hline
  \end{tabular}

  {\footnotesize Fonte: Elaboração própria.}
\end{table}

A escolha de modelos distintos para geração e avaliação foi intencional. O guardrail de fidelidade utiliza um modelo mais robusto do que o modelo responsável pela simplificação, pois sua função é comparar a saída gerada com o material original e identificar possíveis desvios de conteúdo. No contexto deste trabalho, essa escolha foi adotada como uma decisão de projeto para aumentar o rigor da etapa de verificação, buscando reduzir o risco de aprovação de simplificações que apresentem alterações sutis de sentido. Ainda assim, essa verificação não é tratada como garantia absoluta de correção, mas como um mecanismo adicional de controle dentro do fluxo proposto.

No frontend, foram utilizados React 18 e Vite. O React foi adotado por favorecer a construção da interface em componentes reutilizáveis, o que se mostrou adequado para organizar elementos como seleção de perfis, entrada de conteúdo, visualização de resultados, glossário, perguntas frequentes e chat contextual. O Vite foi utilizado como ferramenta de desenvolvimento e build por oferecer configuração simples e ciclo de desenvolvimento rápido para aplicações React. A interface se comunica com o backend por chamadas HTTP e não conhece diretamente os detalhes dos provedores de LLM, dos prompts ou dos guardrails. Essa separação reforça o papel do backend como camada responsável pela lógica da solução.

\section{Representação do perfil do paciente}
\label{sec:representacao-perfil-paciente}

Os perfis de paciente utilizados na implementação são simulados e não correspondem a dados reais. Essa escolha delimita o escopo acadêmico da solução e evita o uso de informações sensíveis de pacientes durante o desenvolvimento. Os perfis servem para demonstrar como diferentes características comunicacionais podem influenciar a forma de apresentação do mesmo conteúdo.

Na implementação, cada perfil contém campos como identificador, nome, idade, escolaridade e área de formação. Esses campos são utilizados como sinais comunicacionais para orientar a adaptação do texto, influenciando aspectos como vocabulário, tom, extensão das frases e nível de detalhamento. Por exemplo, um perfil com menor escolaridade pode receber explicações com frases mais curtas e menor uso de termos técnicos, enquanto um perfil com maior familiaridade com linguagem acadêmica ou da área da saúde pode receber uma versão menos simplificada. Assim, o perfil atua sobre a forma de apresentação do conteúdo, sem modificar sua orientação clínica.

A implementação considera três perfis simulados: Ana, 30 anos, ensino superior e área da saúde; Carlos, 45 anos, ensino médio e área de negócios; e Maria, 68 anos, ensino fundamental e outra área de formação. A elaboração dessas personas contou com o auxílio do ChatGPT para a definição de características comunicacionais coerentes com cada combinação de atributos. Esses perfis foram definidos para representar diferentes níveis de familiaridade com linguagem técnica e diferentes necessidades de adaptação comunicacional. O perfil de Ana representa uma usuária com maior familiaridade com termos de saúde, o perfil de Carlos representa um usuário intermediário e o perfil de Maria representa uma usuária com maior necessidade de linguagem simples, direta e acessível.

A seleção de múltiplos perfis na interface reforça essa proposta. O mesmo conteúdo pode ser processado para diferentes perfis em paralelo, permitindo comparar como a linguagem é adaptada em cada caso. Essa comparação é útil tanto para demonstrar a funcionalidade quanto para apoiar a avaliação da solução. Essa variação entre perfis também permite observar diferentes níveis de simplificação conforme o contexto comunicacional informado, sem que isso signifique personalização clínica da conduta.

\section{Construção dos prompts e integração com LLMs}
\label{sec:construcao-prompts-integracao-llms}

A construção dos prompts foi isolada em funções específicas, concentrando em uma camada própria as instruções enviadas à LLM. O prompt de simplificação define o modelo como uma ferramenta de apoio comunicacional em saúde, apresenta o texto original como fonte principal, descreve o perfil comunicacional do paciente, estabelece regras de simplificação e delimita o formato esperado da resposta. O perfil do paciente é convertido em instruções textuais, orientando aspectos como vocabulário, tom, extensão das frases e nível de detalhamento, sem ser utilizado como base para decisões clínicas.

Essa conversão ocorre de forma programática: a função construtora do prompt seleciona blocos de instrução distintos conforme os atributos do perfil. O nível de escolaridade determina as orientações de vocabulário e complexidade textual; a faixa etária determina o tom sugerido; e a presença de comorbidades adiciona blocos específicos para destacar informações relevantes ao grupo — sempre condicionados ao que está presente no texto original. Dessa forma, o prompt não é um template fixo, mas um documento composto dinamicamente a partir do perfil informado.

Além de orientar a forma da resposta, o prompt também estabelece restrições explícitas para preservar o conteúdo original. Entre elas estão não utilizar conhecimento externo, não alterar dose, frequência, duração ou modo de uso, não suavizar riscos e não criar recomendações médicas novas. Dessa forma, o prompt atua tanto na adaptação comunicacional quanto na delimitação do comportamento esperado da LLM.

Além da simplificação, o prompt solicita a geração de glossário. Nesse caso, a LLM deve identificar termos técnicos presentes no conteúdo e explicá-los em linguagem simples, sem acrescentar informações que não possam ser derivadas do texto original. A resposta é estruturada para permitir que o backend separe o texto simplificado do glossário e dos demais elementos exibidos na interface.

A integração com os modelos de linguagem foi implementada por meio de uma interface de provedor. Com isso, o fluxo principal da aplicação não depende diretamente de uma API específica: o use case de simplificação solicita a geração de texto a partir de um contrato comum, enquanto implementações concretas ficam responsáveis pela comunicação com provedores como Gemini ou outros modelos configurados. Essa abstração reduz o acoplamento e permite trocar ou ajustar o provedor de LLM sem alterar a lógica principal da aplicação.

A implementação também separa conceitualmente o provedor gerador do provedor avaliador. O primeiro é responsável por gerar a simplificação, enquanto o segundo pode ser usado nos processos de avaliação de fidelidade e guardrail. Mesmo que ambos apontem para o mesmo provedor em determinada configuração, essa separação permite usar modelos distintos em cada etapa.

A construção dos prompts e a integração com LLMs, portanto, desempenham papel central na implementação. Elas traduzem os objetivos conceituais da proposta - simplificar, adaptar e preservar o significado - em instruções operacionais e chamadas controladas ao modelo.

A Figura~\ref{fig:prompt-exemplo-simplificacao} apresenta um trecho representativo do prompt enviado à LLM simplificadora. O prompt completo inclui instruções adicionais de preservação, adaptação ao perfil, geração de glossário e formato de resposta. Por ser extenso, o exemplo foi reduzido para evidenciar apenas sua estrutura principal.

\input{figuras/implementacao/prompt-simplify-exemplo.tex}

Além do prompt de simplificação, a implementação utiliza um prompt específico para o guardrail de fidelidade. Esse prompt orienta uma LLM avaliadora a comparar o texto original com a versão simplificada, verificando se as informações incluídas permanecem apoiadas no material de origem. A Figura~\ref{fig:prompt-exemplo-guardrail} apresenta um trecho representativo desse prompt.

\input{figuras/implementacao/prompt-guardrail-exemplo.tex}

\section{Indicador de legibilidade}
\label{sec:indicador-legibilidade}

A implementação calcula o Flesch-PT para o texto original e para a versão simplificada gerada pela LLM. Como apresentado na Seção~\ref{sec:legibilidade-textual}, esse indicador considera o número médio de palavras por sentença e o número médio de sílabas por palavra, mas, na aplicação, esses valores são utilizados apenas como componentes internos do cálculo. O resultado final exibido e considerado na avaliação é o índice Flesch-PT.

No backend, o cálculo é realizado a partir da divisão do texto em sentenças e palavras. As sentenças são identificadas por pontuação final (`.', `!' ou `?') seguida de espaço em branco, descartando fragmentos muito curtos resultantes da divisão. As palavras são obtidas pela divisão em espaços, com remoção de caracteres não-alfabéticos. A contagem de sílabas é feita com a biblioteca \texttt{stress-pt}, que retorna a separação silábica de cada palavra no formato \texttt{pa|la|vra}; a implementação conta o número de segmentos para obter a quantidade de sílabas. A partir desses valores, a aplicação calcula o índice Flesch-PT, limita o resultado ao intervalo de 0 a 100 e associa a pontuação obtida à respectiva classificação de legibilidade.

Os valores calculados para o texto original e para o texto simplificado são retornados pelo backend junto ao resultado da simplificação e exibidos na interface, permitindo comparar as duas versões. O indicador é utilizado apenas como medida automática de legibilidade textual, não como avaliação de precisão médica, segurança clínica ou qualidade semântica da resposta. Dessa forma, um aumento na pontuação do Flesch-PT indica sinais de maior facilidade de leitura, mas não garante, por si só, que a simplificação preservou corretamente o sentido do conteúdo original.

\section{Interface da aplicação}
\label{sec:interface-aplicacao}

A interface da SimplyHealth foi construída para demonstrar o fluxo principal da solução e permitir a comparação entre versões geradas para diferentes perfis. Ela pode ser entendida em duas etapas principais: a tela inicial e a tela de resultados.

Na tela inicial, ilustrada na Figura~\ref{fig:tela-inicial}, o usuário seleciona um ou mais perfis simulados de pacientes, insere texto livre ou envia um arquivo PDF/TXT e escolhe parâmetros de processamento. Quando aplicável, também pode indicar se imagens presentes nos PDFs devem ser consideradas. Essa etapa concentra as decisões necessárias para iniciar o processamento.

A possibilidade de selecionar múltiplos perfis é relevante para o objetivo do trabalho. Ela permite observar como o mesmo texto técnico pode ser adaptado de formas diferentes conforme o perfil comunicacional escolhido. Assim, a interface não apenas executa a simplificação, mas também torna visível a lógica de adaptação ao perfil.

\begin{figure}[htbp]
  \centering
  \includegraphics[width=\textwidth]{figuras/implementacao/tela-inicial.png}
  \caption{Tela inicial da SimplyHealth: seleção de perfis, inserção de conteúdo e parâmetros de processamento.}
  \label{fig:tela-inicial}
  {\small Fonte: elaboração própria.}
\end{figure}

Na tela de resultados, a aplicação apresenta os resultados em cartões por paciente. Cada cartão exibe o texto simplificado e o indicador de legibilidade, conforme ilustrado na Figura~\ref{fig:tela-resultados-cima}. A área inferior do mesmo cartão, mostrada na Figura~\ref{fig:tela-resultados-baixo}, apresenta o glossário recolhível, perguntas sugeridas e acesso ao chat contextual. Também há uma área para consulta do material original, permitindo comparar a fonte com a versão simplificada.

\begin{figure}[htbp]
  \centering
  \includegraphics[width=\textwidth]{figuras/implementacao/tela-resultados-cima.png}
  \caption{Tela de resultados da SimplyHealth — área superior: texto simplificado e indicador de legibilidade por perfil.}
  \label{fig:tela-resultados-cima}
  {\small Fonte: elaboração própria.}
\end{figure}

\begin{figure}[htbp]
  \centering
  \includegraphics[width=\textwidth]{figuras/implementacao/tela-resultados-baixo.png}
  \caption{Tela de resultados da SimplyHealth — área inferior: glossário recolhível, perguntas frequentes e chat contextual.}
  \label{fig:tela-resultados-baixo}
  {\small Fonte: elaboração própria.}
\end{figure}

Essa organização favorece a avaliação da solução, pois permite observar simultaneamente o conteúdo original, o texto simplificado e os elementos de apoio à compreensão. Além disso, a presença de resultados separados por perfil facilita a análise das diferenças entre versões, aspecto importante para verificar se a adaptação comunicacional está ocorrendo de forma coerente.

A implementação descrita neste capítulo materializa a proposta da SimplyHealth em uma aplicação funcional, composta por frontend, backend, integração com LLMs, mecanismos de guardrail, recursos de apoio à compreensão e indicador de legibilidade. A partir dessa implementação, torna-se possível avaliar a solução em relação ao funcionamento do fluxo proposto, à atuação dos guardrails e à qualidade das simplificações geradas. O próximo capítulo apresenta a metodologia de avaliação da solução desenvolvida.
